\section{怎样修改病句}

我们在中学生的作文里，有时会碰到这样一类的句子：

（1）今天，我们是在党的阳光照耀下健康成长的青年是多么幸福！

（2）是否学好了语文，标志着为学习其他各门功课打下坚实的基础。

（3）严老师一走上讲台，就用猫头鹰似的目光看了大家一眼，喧闹的教室突然鸦雀无声了。

我们看，第一句把“我们是……青年”和“青年是多么幸福”两个句子的意思杂糅在一起，使它既不是单句，又不是复句，也就是说，它不像个句子，弄得读者莫名其妙。第二句，前面用了既肯定又否定的两面词“是否”，后面却用只表示肯定的一面词“标志着”与之呼应，而我们从常识中可以知道，“是”是能够“标志着”打下了坚实的基础，“否”却是不能作为打下坚实基础中的“标志”的。第三句，形容一位教风严肃又富有威信的好老师的目光。不用带褒意感情色彩的词语，却用带贬意感情色彩的比喻词语“猫头鹰似的”，显然不恰当了。

人们把不能准确、鲜明地表达作者思想感情的句子叫做病句。这一类有毛病的句子，由于表达思想感情不准确或者不鲜明，别人就不能理解甚至会误解句子要表达的意思。这样就达不到交流思想感情的目的，因此，这类病句必须加以修改。

那么，怎样才能修改好病句呢？

\textbf{1. 弄清病因是修改病句的前提。}

毛泽东同志说过：“一个合逻辑，一个合文法，一个较好的修辞，这三点请你们在写文章的时候注意。”而病句的出现正由于对这三点不注意，它们总是有违反逻辑、违反语法规则或违反修辞一般要求的地方。象前面所举的第一句，就违反语法规则，把两种从不同角度说的意思硬凑在一起，造成结构和意思上的混乱，使它既不成为单位句，也不成为复句。找出这个病因，我们就可以删去第一个“是”字，使全句成为结构完整、意思分明的单位；或者在“青年”后面加上逗号和“我们”二字，使全句成为由两个意义紧密相关的分句组合而成的复句。再如前面所举的第二句，就违反逻辑，前言不对后语，因为前面既有肯定又有否定，而后面只有肯定，这使前后失去了呼应。找到了病因，我们就可以对症下药：或者删去开头的两面性词语“是否”，使全句前后呼应，保持统一的肯定角度；或者在“标志着”后面，加上“是否”二字，也能使全句前后呼应，同时保持了统一的既有所肯定又有所否定的角度，这样就符合事理了。又如第三个句子，就违反了修辞的一般要求，因为在第一个句子中感情色彩必须与句意完全一致。如果在一个该用褒义词语的句子中用了贬义词语，或者在该用贬义词语的句子中用了褒义词语，都不能正确地表达思想感情。找到了这句话的病因，就可以把“猫头鹰似的”改成“威严的”，全句的感情色彩就一致了。

当然，病句的病因是多种多样的，对于具体的病句要进行具体的分析，但是总起来说，不外乎不合语法、不合逻辑和修辞不当三种类型。我们面对一个病句，首先应从语法方面检查病因。如果属于语法方面的毛病，就再具体分析它属于语法的哪一方面的毛病，以便着手修改。如果语法方面没有毛病，就应从是否不合事理、违反逻辑方面考虑了。比如：“凡在科学研究上有杰出成就的人，都是在物质条件十分艰苦的情况下，经过顽强刻苦的努力才获得成功的。”这句话在语法方面并没有毛病，但是从逻辑上检查，它却存在着以偏概全、判断不当的错误，因此只有把“凡”字改为“许多”（或者“不少”），并且删去“都”字，这个句子才是周密无误的。至于修辞方面的毛病，不外乎用词色彩不恰当，比喻不恰当，夸张不恰当，拟人不恰当，等等，只要用心体会，仔细辨析，是不难发现它的病因的。

\textbf{2. 掌握病句类型是修改病句的基础。}

修改病句不仅要弄清此句是语法上有毛病，还是逻辑上有毛病，还是修辞上有毛病，还是有几个方面的毛病，而且要弄清它属于哪一种具体的病句类型，只有这样，才能准确判断这个病句的毛病究竟在什么地方，是一个地方有毛病还是几个地方有毛病。有了这个基础，着手修改就比较好办了。

病句有哪些基本类型呢？我们仅就常见的语法错误、逻辑错误和修辞错误，列为（1）、（2）、（3）三表，读者可以细心体会，并且以此类推，以求触类旁通。
\newpage
\textbf{(1) 犯语法错误的常见病句类型简表}

\noindent
\begin{tabular}{|lp{5em}|p{10em}|p{4em}|}
	\hline
	\multicolumn{2}{|l|}{错误类型} & 病句举例 & 订正\\ \hline
	\multicolumn{1}{|l|}{\multirow{3}{*}{用词错误}} & 生造词语  & 祖国的山河是多么伟丽啊！      & 雄伟壮丽 \\ \cline{2-4} 
	\multicolumn{1}{|l|}{} & 错用同义词 & 我们要象保卫眼珠一样地保卫秧苗。 & 保护\\ \cline{2-4} 
	\multicolumn{1}{|l|}{} & 错用名词 & 这个城市如今显得更青春、更美丽了。 & 年青\\ \cline{2-4}
	\multicolumn{1}{|l|}{} & 错用动词 & 我们怀着瞻仰的心情踏上纪念碑的台阶。 & 崇敬\\ \cline{2-4}
	\multicolumn{1}{|l|}{} & 错用形容词 & 我们必须发达自己的工业和农业。 & 发展\\ \cline{2-4}
	\multicolumn{1}{|l|}{} & 错用数量词 & 现在我们用的无线电望远镜有二样类型。 & 两种\\ \cline{2-4}
	\multicolumn{1}{|l|}{} & 错用代词 & 要是有困难，咱们给你想办法。 & 我们\\ \cline{2-4}
	\multicolumn{1}{|l|}{} & 错用副词 & 对于那件事，我记得非常清清楚楚。 & 删去“非常”\\ \cline{2-4}
	\multicolumn{1}{|l|}{} & 错用介词 & 我们应该把马列主义理论武装自己的头脑。 & 用\\ \cline{2-4}
	\multicolumn{1}{|l|}{} & 错用助词 & 张敏正在细心地配了第二服药。 & 着\\ \cline{2-4}
	\multicolumn{1}{|l|}{} & 错用连词 & 我班同学都是本市新海街或新中街的人。 & 和\\
	\hline
\end{tabular}

\noindent
\begin{tabular}{|lp{5em}|p{8em}|p{6em}|}
	\hline
	\multicolumn{2}{|l|}{错误类型} & 病句举例 & 订正\\
	\hline
	\multicolumn{1}{|l|}{\multirow{3}{*}{造句错误}} & 缺主语成分  & 通过学习，使我提高了认识。 &删去“使” \\ \cline{2-4} 
	\multicolumn{1}{|l|}{} & 缺谓语成分 & 今天下午我们二班跟三班拔河友谊赛。 & 前面加“进行”\\ \cline{2-4}
	\multicolumn{1}{|l|}{} & 缺宾语成分 & 我们要学习白求恩同志毫不利己专门利人。 & 后面加“的精神”\\ \cline{2-4}
	\multicolumn{1}{|l|}{} & 主谓搭配不当 & 革命的队伍比以前扩大和提高了。 & 删去三字\\ \cline{2-4}
	\multicolumn{1}{|l|}{} & 动宾搭配不当 & 我们要节省一切不必要的开支和浪费。 & 删去三字\\ \cline{2-4}
	\multicolumn{1}{|l|}{} & \multirow{2}{*}{定、状、补和中心词搭配不当} & 革命导师都有非常渊博的革命知识和经 & \\ \cline{3-4} 
	\multicolumn{1}{|l|}{} & & 王民同学总是热火朝天地学习。 & \\ \cline{3-4}
	\multicolumn{1}{|l|}{} & & 我们把教室打扫得干干净净，整整齐齐了。 & 前面加“把桌凳安排得” \\
	\hline
\end{tabular}

\noindent
\begin{tabular}{|lp{3em}|p{8em}|p{8em}|}
	\hline
	\multicolumn{2}{|l|}{错误类型} & 病句举例 & 订正\\
	\hline
	\multicolumn{1}{|l|}{\multirow{3}{*}{造句错误}} & 词序不当 &  我不禁想起了许多王老师的动人事迹。 &  移到“动人”前\\ \cline{2-4}
	\multicolumn{1}{|l|}{} & 句子成分多余 & 大家听了非常高兴得很。 & 保留一个\\ \cline{2-4}
	\multicolumn{1}{|l|}{} & 句子结构混乱 & 这件事使我有很大的教育意义。 & 改“使”为“对”，或在“使我”后面加上“很感动，对我”一类文字\\ \cline{2-4}
	\multicolumn{1}{|l|}{} & 复句关系不对 & 如果破除了迷信，改革了技术，那才不但能大大加快生产的进度，而且能大大提高产品的质量。 &  “如果”改为“只有”；或“才”改为“就”\\
	\hline
\end{tabular}

\vspace{2em}
\textbf{(2) 犯逻辑错误的常见病句类型简表}

\noindent
\begin{tabular}{|ll|p{10em}|p{6em}|}
	\hline
	\multicolumn{2}{|l|}{错误类型}                         & 病句举例                 & 订正      \\ \hline
	\multicolumn{1}{|l|}{\multirow{4}{*}{概念不清}} & 偷换概念 & 画廊中的画，除了王老师外，都是学生画的。 & 后面加“画的” \\ \cline{2-4} 
	\multicolumn{1}{|l|}{}                      & 种属混乱 & 李小波在本学期参加了共青团员。      & 删去“员”   \\ \cline{2-4} 
	\multicolumn{1}{|l|}{}                      & 相互排斥 & 开演前五分钟，观众们争先恐后地鱼贯入场。 & 删去“鱼贯”  \\ \cline{2-4} 
	\multicolumn{1}{|l|}{}                      & 并列不当 & 饲养场的家禽、鸡、鸭、鹅数量猛增。    & 删去“鱼贯”  \\ \hline
\end{tabular}

\noindent
\begin{tabular}{|lp{5em}|p{10em}|p{5em}|}
	\hline
	\multicolumn{2}{|l|}{错误类型}                            & 病句举例                 & 订正         \\ \hline
	\multicolumn{1}{|l|}{\multirow{6}{*}{判断错误}} & 自相矛盾    & 在团支委会上，大家互相作了自我批评。   & “互相”改为“各自” \\ \cline{2-4} 
	\multicolumn{1}{|l|}{}                      & 主、宾配合不当 & 欧阳珠的故乡是江苏宜兴人。        & 删去“人”      \\ \cline{2-4} 
	\multicolumn{1}{|l|}{}                      & 前后不能呼应  & 他紧张地站起来，盼望老师能问较难的问题。 & 改为“浅”      \\ \cline{2-4} 
	\multicolumn{1}{|l|}{}                      & 正反颠倒    & 如今出言不逊的现象大大进步了。      & 改为“减少”     \\ \cline{2-4} 
	\multicolumn{1}{|l|}{}                      & 主客颠倒    & 故乡发展的情况对于我是十分熟悉的。    & 二者应对调      \\ \cline{2-4} 
	\multicolumn{1}{|l|}{}                      & 乱用多次否定  & 不论谁，工作中难免没有这样那样的错误。  & 删去“没”      \\ \hline
\end{tabular}

\vspace{2em}
\noindent
\begin{tabular}{|lp{5em}|p{10em}|p{5em}|}
	\hline
	\multicolumn{2}{|l|}{错误类型}                               & 病句举例                   & 订正          \\ \hline
	\multicolumn{1}{|l|}{\multirow{4}{*}{推理荒谬}} & 根据不正确      & 二班的赵杰是个阴谋家，他的话不能全信。    & 改为“常撒谎”     \\ \cline{2-4} 
	\multicolumn{1}{|l|}{}                      & 理由不充足      & 赵厂长是个称职的干部，因为他工龄长。     & 加上“能力强，觉悟高” \\ \cline{2-4} 
	\multicolumn{1}{|l|}{}                      & 强加因果关系     & 这学期我认真学习了语文，因而数学学得不够好。 & 改为“但是”      \\ \cline{2-4} 
	\multicolumn{1}{|l|}{}                      & 轻率概括（以偏概全） & 李华、李平和李红成绩出众，所以姓李的都很好。 & 前面加“我班”     \\ \hline
\end{tabular}

\textbf{(3) 犯修辞错误的常见病句类型简表}

\noindent
\begin{tabular}{|l|p{12em}|p{6em}|}
	\hline
	错误类型  & 病句举例                          & 订正        \\ \hline
	乱用贬义词 & 要是没有共产党，我们一家早就完蛋了             & 删去“蛋”     \\ \hline
	乱用褒义词 & 来犯的敌人大半被歼灭，剩下的几个非常灵活，连忙钻进了山洞。 & 改为“狡猾”    \\ \hline
	乱用形容词 & 他红光满面，如火如荼地说：“谢谢大家的批评，我一定改。”  & 改为“面红耳赤”  \\ \hline
	比喻不当  & 盛夏的中午，太阳似明镜一样悬在当空，烤得石头都烫人。    & 改为“巨大的火球” \\ \hline
	比喻失实  & 他头上的汗珠象铅球似地往下落。               & 改为“黄豆”    \\ \hline
	夸张不当  & 棉桃一个泰山大，十辆汽车往社搬。              & 改得能使前后搭配  \\ \hline
	拟人不当  & 晚上，只听见几只蟋蟀在窗下大喊大叫。            & 改为“轻轻地歌唱” \\ \hline
	拟物不当  & 我刚骂了弟弟一句，妈妈就张牙舞爪批评我。          & 改为“不容分说”  \\ \hline
\end{tabular}

\vspace{2em}
\textbf{3. 找准具体毛病是修改病句的关键。}

一个病句摆在我们面前，我们可以根据常见病句类型弄清这个病句，是什么性质的病句，犯的是语法上的、逻辑上的还是修辞上的毛病，然后具体分析病句的每个成分以及这些成分之间的关系，找准具体毛病出在哪儿究竟是一个地方有毛病还是几个地方有毛病。只要这一步做好了，怎样订正这个病句也就好办了。例如：

（1）这时，全场所有人的眼睛都集中到了大会主席台上。

（2）教育战线和大专院校是培养社会主义建设人才的“一支队伍”。

（3）一直到五月下旬，彗星的亮度仍然在肉眼能见到的限度之上。

（4）在那次激烈的战斗中，我们击落了敌人大批的飞机和坦克。

（5）晴空万里，春风拂面，柳枝轻舞，青烟直上。

第一句的主语“眼睛”和谓语“集中”不能配搭，从逻辑方面看，它把“眼睛”和“眼光”两个概念混淆起来了。

第二句的“教育战线”与“大专院校”不是相应相称的概念，不能并列在一起；同时，整句的判断犯了主宾不能配合的毛病，因为“教育战线和大专院校”不能是“一支队伍”。

第三句“在······限度之上”，与“亮度仍然······”是互相矛盾的，因为在限度“之上”是说超过了限度，就不会“仍然”“能见”了。

第四句的谓语“击落”与宾语“坦克”不能配搭，因为“飞机”可以“击落”，“坦克”是不能“击落”的。

第五句的“直上”“春风拂面”是互相矛盾的，因为既然春风能够“拂面”，柳枝能够“轻舞”，青烟就不可能是“直上”了。

当我们找准了具体毛病之后，就可以着手修改了。

\textbf{4.尽量少改和不损害原意是修改病句的基本要求。}

一般说来，作修改练习用的病句并不是毛病百出的，而是只要改动一两处就可以通顺的，因此我们在找准毛病所在之处，稍加订正就行了。改动的方法包括删去多余的字词，增加必要的字词，更换某些字词，或者调整某些字词的次序等。不管采用何种改法，基本要求是一样的，那就是以改得很少而又能保持原意的改法为好。

在具体修改过程中，应该注意：不要乱涂乱改，或者撇开原句另造新句，因为这样做不符合修改病句的基本要求，而且可能使有毛病的地方没得到改正，反倒把没有毛病的地方改出毛病来，这是一；修改时，要在尽可能保持原意的前提下，力求只改动一处就使句子通顺，同时不要把原句中有数的表意成分一笔勾销，从而损害了原意，就象医生看病一样，不能因为病人手指上生个小疮而把整个手指截掉，这是二；在有两种以上改法的情况下，需要用心体会原句的句意，比较这两种以上的改法，寻求改动既少又能保持原意的最佳改法，这是三。例如：

上面第三点中所举的第一个病句，当找准它的具体毛病是在主语“眼睛”和谓语“集中”不能配搭，就可以考虑两种改法：一是保留主语“眼睛”而把谓语“集中到了”改成“盯着”；二是保留谓语“集中”而把主语“眼睛”改成“眼光”。比较这两种改法，可以看出第二种改法改动既少，又最能传神地表达原意，因此就选择这第二种改法。

上面第三点中所举的第二个病句，当我们找准了它的具体毛病后，发现可能有四种改法：一是删去“教育战线和”，二是删去“和大专院校”；三是改动“教育战线和大专院校”；四是改动“一支队伍”。总观原句句意，“教育战线和大专院校”都是指教育界的事，是有效的表意成分，不能都删去，因此第三种改法不能成立；“教育战线”比“大专院校”范围大，包括了“大专院校”，在这个句子中删去“大专院校”并没有改变基本句意，因此不采取第一种改法而采取第二种改法；但是在采取了第二种改法后，主语“教育战线”与宾语“一支队伍”不能配搭，因此可以把“一支队伍”再改为“一条重要战线”，全句就通顺了。

上面所举的第三个病句，毛病既然在于前后矛盾，我们只要将“在……限度之上”的“上”字改为“内”或“中”就行了，如果要把“仍然在”的“在”字改为“没有超过”，句意也通顺，只是改动较多，不如第一种改法好。

上面所举的第四个病句，既可以删去“和坦克”三字，也可以将“击落”改为“击毁”。两种改法都能使句子通顺，但是“坦克”是该句的有效表意成分，不应删掉，所以选择改动“击落”为好，不仅改动少，又能保持原意。

上面所举的第五个病句，毛病也出在前后矛盾上。那么，我们只要将“青烟直上”的“直上”改为“缭绕”，就能做到句意圆合了。当然，保留“青烟直上”，而把“春风拂面、柳枝轻舞”改成“春风不起，柳枝低垂”，句子也是通顺的，但是改动既多，又与句子的原意有出入，我们就应该放弃这种改法了。

修改病句是一项培养遣词造句能力的有效训练方式，它需要我们切实掌握语法、逻辑、修辞的有关知识，并能综合地灵活地运用。我们只要记住各种病句的类型，懂得辨析和修改病句的方法，经过一番修改病句的实践，这个问题也是不难解决的。
\newpage
